\section{怎样朗读文章}

同一篇文章，给善于朗读和不善于朗读的人读起来，可以产生不同的效果。这说明朗读也是一种重要的语文能力。我们通过朗读，不仅能把作品的思想内容准确、鲜明、生动地表达出来，加深了对文章的精神实质的理解，而且也有助于掌握运用语言的规律，丰富语汇，提高语文水平和表达能力，甚至还有利于纠正方言，更好地学习和推广普通话。

那么，怎样才能朗读得好呢？

从思想方面来说，我们首先要深入分析文章的思想内容，深刻理解文章的精神实质，把握住文章的思想基调和感情起伏，然后以满腔的政治热情和鲜明的爱憎感情把文章的思想感情恰如其分地传达出来。例如朗读《谁是最可爱的人》，我们就应该象作者魏巍一样充满了对志愿军战士的火热的爱，用喷薄而出的革命激情去歌颂革命战士那种爱祖国、爱人民、恨敌人的崇高思想和革命英雄主义、爱国主义、国际主义的伟大精神，从而使自己沉浸在对“最可爱的人”的无比崇敬、无比爱戴的感情波涛之中，并激起听者同样的感情。又如朗读《包身工》，我们在掌握本文对包身工悲惨遭遇的同情和对这种罪恶制度的抨击这个思想感情的基调之后，就应该怀着作者夏衍那样鲜明、强烈的爱憎，用沉重、愤懑的音调，把文章的思想倾向传达出来，使自己受到教育，也使听者受到感染。

从语言方面来说，我们首先要读得正确，然后还得注意停顿、重音、语调等朗读常识和基本技巧，这才有可能收到抑扬顿挫、通感荡气的朗读效果。所谓读得正确，包括普通话要说得准确、流利，避免漏字、添字和颠倒词语，还得吐字清楚，不含糊其词。这尽管是读好文章的起码条件，但是却不可忽视，因为这是朗读的基本功。试想，如果普通话说得不伦不类，整里整扭，怎能期望传情达意，感染听众？如果掉字、添字、词语颠倒，使人听了莫名其妙，又怎能表达文章的思想感情？至于咬字不清，话在嗓子里窝困着，使人连意思都没听出来，当然更不可能打动听众了。

读得正确是朗读的基本要求，要朗读得好，还得了解一些朗读的基本常识和朗读技巧才行。

\textbf{1. 掌握恰当的停顿。}

停顿是指朗读时或长或短的间歇。句子有长有短，短句子可以一口气读完，长句子就要在适当的地方作些停顿。停顿固然是生理上的需要，但更主要的是为了清楚而准确地表达语意，传达感情，好让听者有时间来思考、理解朗读的内容，以引起共鸣。

朗读时，停顿时间的长短一般可根据标点符号来确定。顿号停顿最短，逗号停顿稍长，分号、冒号又比逗号的停顿稍长，而句号（包括问号、感叹号，它们都有表示句子结束的作用）的停顿则比分号、冒号又稍长。注意到这一点，我们对一句话的朗读如何掌握停顿的时间就大致合乎要求了。

句与句之间有个较长的停顿，段与段之间则需要更长的停顿。这是因为一段话完了之后，要转移或转折到另一段意思上去，必须有一个明显的标志。因此，朗读时，我们应该使听众能听出句与段来，而不应该像打连珠炮一样地一口气读到底，读得连气都喘不过来，哪里还有时间让听众思索、回味所读的思想内容呢？

一般说来，段与段之间有个较长的停顿，可以收到“此时无声胜有声”的效果；不过，这也不能绝对化，还得根据内容表达的要求来处理某些特殊的情况。例如《为了六十一个阶级弟兄》里，写为了抢救六十一个阶级弟兄的生命，北京的空军战士将一箱情意比泰山还重的特效药连夜空投到平陆之后：

“盼！盼！----在张村公社医院的大门口，社员们、医护人员正焦急地盼望着······

汽车开来了！----好！

马上拿下药箱！

马上注射！”

朗读时，这里的停顿，就不能一般地根据标点符号来处理，也不能将段与段之间留下更长时间的间歇，而应该用较短的停顿来读，这样才能渲染出当时万分紧急的情势和人们惊喜交集的深厚的阶级感情。

另外，有些表意严密，并列成分或附加成分较多的长句子，尽管句子中没有用标点符号，但是为了把意义表达得更清楚，我们可以根据句子所表现的意义和所强调的词语，在朗读时可以适当地运用停顿。例如：

“马克思列宁主义的普遍真理 $\vert$ 
一经和中国革命的具体实践相结合 $\Vert$ 
就使中国革命的面目 $\vert$ 
为之一新。$\Vert$”
（用$\vert$隔开的地方，表示根据句子的结构和意义可以稍作停顿；用$\Vert$隔开的地方，表示停顿较长。）

“城北丹崖山峭壁上$\vert$ 
那座凌空欲飞的蓬莱阁$\Vert$ 
尤其有气势。$\Vert$ ”

（用 $\vert$ 
隔开的地方，可以稍作延长，表示句中的自然停顿；用 $\Vert$ 
隔开的地方，表示强调主语或句意结束，停顿可以稍长一点。）

在朗读诗歌的时候，不但要注意停顿，还要注意节拍。节拍之间可以是短暂的停顿，也可以是字音的延长。诗歌的语言富有音乐美，按节拍朗读，可以有鲜明的节奏感；同时诗歌的语言精炼、含蓄，按节拍朗读，也便于听者思考、理解和体味。例如：（---表示延长）

羊群 $\Vert$ 走路 --- 靠头羊--- $\Vert$，

陕北 $\Vert$ 起了 --- 共产党 --- $\vert$。

千里的 $\Vert$ 雷声 --- 万里的 $\vert$ 闪 --- $\vert$，

陕北 $\Vert$ 红了 --- 半个天 --- $\vert$。
\begin{flushright}
——李季《王贵与李香香》
\end{flushright}

远远的 $\vert$ 街灯 $\vert$ 明了 --- $\vert$，

好象是 --- 闪着 $\vert$ 无数的 $\vert$ 明星 --- $\Vert$。

天上的 $\vert$ 明星 $\vert$ 现了 --- $\vert$，

好象是 --- 点着 $\vert$ 无数的 $\vert$ 街灯 --- $\Vert$。
\begin{flushright}
——郭沫若《天上的街市》
\end{flushright}

\textbf{2. 掌握必要的重音}

朗读时，一个句子中的某些词语要读得重些，某些词语要读得轻些。读得重些的词语就是句子的重音。重音对于区分一般语意和必须强调的意义具有特别重要的作用。我们应该掌握它。

重音一般分为语法重音和逻辑重音两种。

(1) 语法重音是指在一般句子中自然而然要念得重些的音节，所以也叫自然重音。比如说，在一些比较短的句子里，谓语往往要读得重些，句中有定语、状语和补语时，它们往往比主语、谓语和宾语读得重一些；有些作比喻的句子，重音总是落在用来作比喻的词语上。例如：

①\ 中国人民站起来了！（重音在谓语上）

②\ 自然是伟大的，然而人类更伟大。（重音在定语和状语上）

③\ 我们要刻苦地钻研，不懈地努力。（重音在状语上）

④\ 大家激动地跳了起来。（重音在补语上）

⑤\ 海燕像黑色的闪电，在高傲地飞翔。（重音在比喻词和状语上）

(2)\ 逻辑重音是指为了突出句子中某个意思而特别加以强调的音节。同样一句话，由于逻辑重音不同，意思会有显著的变化。比如某个中学开家长座谈会，门前打起这样一幅横标：“亲爱的爸爸妈妈欢迎您！”如果逻辑重音读得不同，意思也跟着变了：

①\ 亲爱的爸爸妈妈欢迎您！（强调“欢迎”，不是不欢迎。这是横标的本意。）

②\ 亲爱的爸爸妈妈欢迎您！（强调爸妈欢迎的对象是“您”，而不是别的什么人。与本意大相径庭。）

③\ 亲爱的爸爸妈妈欢迎您！（强调欢迎“爸爸妈妈”，而不是爷爷奶奶什么人。与本意有出入。）

④\ 亲爱的爸爸妈妈欢迎您！（强调欢迎“亲爱的”，而不欢迎不亲爱的。与本意有出入。）

[按：若是再加上在不同地方的停顿，句意就更有悬殊。如“亲爱的爸爸妈妈，欢迎您！”“亲爱的爸爸，妈妈欢迎您！”“亲爱的，爸爸妈妈欢迎您！”这真是差之毫厘，谬以千里了。]

要确定文章中某个句子的逻辑重音，必须从分析文章的思想内容着手，把前后句联系起来研究。例如：

前年的今日，我避在客栈里，他们却是走向刑场了；去年的今天，我在炮声中逃在英租界，他们则早已埋不知哪里的地下了；今年的今日，我才坐在旧寓里，人们都睡觉了，连我的女人和孩子。我又沉重的感到我失掉了很好的朋友，中国失掉了很好的青年，我在悲愤中沉静下去了，不料积习又从沉静中抬起头来，写了以上那些字。
\begin{flushright}
	--- 鲁迅《为了忘却的记念》
\end{flushright}

鲁迅先生这段话的基调是“悲愤”，既表现了对国民党反动派迫害自己、屠杀革命战友的无比愤怒和仇恨，又表现了对那些为中国人民的解放而英勇就义的革命战友的无限哀悼和怀念。在行文中，还采取了对照的手法，来突出自己对牺牲了的革命战友的敬意和极其沉痛的心情。因此，凡是能突出强调这种“悲愤”感情的词语，都该用逻辑重音来读。

再如：

是谁创造了人类世界？

是我们劳动群众。

一切劳动者所有，

哪能容得寄生虫！
\begin{flushright}
	--- 鲍狄埃《国际歌》
\end{flushright}


“谁”是几千年来唯物史观和唯心史观斗争的焦点；
“劳动群众”是共产主义者的斩钉截铁的回答。为了消灭剥削阶级，根除“劳者不获、获者不劳”这个阶级社会的根本矛盾，必须将一切劳动果实归“劳动者”所有；“哪”这个词正表现了无产阶级消灭剥削阶级、实现共产主义的坚定决心。所以这些词语必须确定为逻辑重音，着重强调地读出来。

表示重音的方法，一是重读——音量增大，声调升高（或者把声音压得低沉），声音拉长；二是在重音词语的前后，作适当的停顿。

\textbf{3. 注意语调的变化。}

语调是指一句话读音的高低、快慢、强弱的不同变化，特别是指高低升降的变化。这种语调的变化决定于表达的意思不同，对事物的态度不同，使用的语气也就不同。

语调变化多端，大致可分为四种：

(1) 升调——语调由低到高。疑问句、反问句末尾的语气上扬，属于升调。例如：

①\ 怎样保卫每一寸土地呢？怎样使每一寸土地都发挥它的巨大潜力，一天天更加美好起来呢？

②\ 中国人失掉自信力了吗？

(2) 降调——语调由高而低。陈述句、祈使句、感叹句一般末尾稍降。例如：

①\ “我真傻，真的，”她说，“我早知道雪天是野兽在深山里没有食吃，会到村里来；我不知道春天也会有。……”
（《祝福》）

②\ “多多头！打死我也不怨你，只求你不要说出来！”
（《春蚕》）

③\ “促进日中人民的友谊，也是斗争的一部分啊！”
（《樱花赞》）

(3) 平调——语调始终平直。用来表示思索的句子，末字拖得较长，属于平调。例如：

①\ 倘有别的意思，又因此发生别的事，则我的答话委实该负若干的责任……。（《祝福》）

②\ 他让自己的思想远远走到前面去：“无数的世界在无穷无尽的宇宙的广阔胸怀中产生、发展、灭亡，又重新产生……”（《火刑》）

(4) 曲调——语调由高而低后又转高。用于讽刺、嘲笑等语句，表示复杂的感情。例如：

我还记得，“国共合作”时代，通信和演说，称赞苏联，是极时髦的，现在可不同了，报章所载，则电杆上写字和“×党”，捕房正在捉得非常起劲，那么，为将自己的论敌指为“拥护苏联”或“×党”，自然也就得合时，或者还会得到主子的“一点恩惠”了。但倘说梁先生意在要得“恩惠”或“金榜”，是冤枉的，决没有这回事，不过想借此助一臂之力，以济其“文艺批评”之穷罢了。所以从“文艺批评”方面来看，就还得在“走狗”之上，加上一个形容字：“乏”。
\begin{flushright}
	---鲁迅《“丧家的”“资本家的乏走狗”》
\end{flushright}

鲁迅先生这段话，嘻笑怒骂，用笔曲折，穷形极相，酣畅淋漓，有力地揭露了梁实秋的阶级本质。我们应该深刻理解作品的思想内容，有自己真切的感受，只有这样，才能用恰当的语调把作品的思想感情准确地表达出来。

当然，朗读不同于演戏。有点表情朗读是可以的，然而切切不可矫揉造作，阴腔怪调，因为那样的朗读失掉了准确传达作品的思想感情的意义，而显得舍本逐末了。
\newpage